% ===== §5 Epistemology: the diagnosis of curvature =====
\section{Epistemology: The Diagnosis of Curvature and the Reversal of the Burden of Proof}
\label{sec:epistemology}

\noindent
The synthesis gave a criterion: the good circle is the one of positive holonomy, the spiral that returns to its root while accruing a remainder. But a criterion of what a thing \emph{is} is not yet a procedure for telling, from outside, \emph{that} it is---and the gap between the two is the whole matter of this section. How, and how far, can the sign of the curvature be read? The section's answer constrains everything that follows, for it arrives at a conclusion that is, at first, disappointing and, on reflection, the very hinge of the paper's practical doctrine: that there is no shortcut to the diagnosis, no external metric that would relieve us of a sustained and fallible attention---and that this absence is not a defect of the theory but a demand of it.

\subsection{The trending-toward-the-good of natural dynamics}
\label{sec:trending}

Begin with an observation that, rightly understood, reverses a burden of proof we scarcely knew we carried. The Daoist tradition offers two homely images, and their homeliness is part of their force, for they record something everyone has seen. A wound left alone will, very often, heal of itself; it is the picking at it, the over-cleaning, the repeated anxious probing, that prevents a self-organizing process wiser than our intervention from doing its work. A body fallen into water will, if it ceases to thrash, float of itself; it is the struggling---the forced action---that defeats the buoyancy by which one would otherwise have risen. \zh{沉疴能自疗，沉溺者不挣则浮}: the deep affliction can heal itself; the one who is drowning, if he does not struggle, floats.

What these images teach is not a vague quietism but something precise about the curvature of certain systems. The wound, the body in water, the convalescing relation: each has a \emph{ground state} that is already a positive curvature, a default tendency toward recovery and self-continuation that operates unless it is interfered with. The body does not need to be taught to float; it floats by its nature, and needs only to stop doing the thing that prevents it. The wound does not need to be made to heal; it heals by its nature, and needs only to be left unpicked. \emb{The default tendency of a natural dynamics is often already toward the good---toward positive holonomy---and such a system goes vicious not for want of intervention but because it has been deformed by intervention.}

This gives \emph{wu wei} a far stronger ground than mere prudence, and the strength of the ground is worth making explicit, for it is easily underestimated. The counsel to ``let be'' is usually heard as a counsel of caution---do less, lest you make things worse---a kind of epistemic humility about our limited foresight. But the claim here is not merely that we are too ignorant to intervene well; it is that the system's own dynamics are, in the relevant cases, already oriented toward the good, so that intervention is not a risky improvement upon a neutral baseline but, characteristically, a deformation of a baseline that was already positive. This reverses a burden of proof:

\begin{claim}[The reversal of the burden of proof]
If the default tendency of natural dynamics is toward the good---if self-continuation and recovery toward positive holonomy are the ground state---then the burden of proof is reversed. It is not \emph{letting be} that requires justification, but \emph{intervening}. The default presumption is that the system will, of itself, tend toward the good; what stands in need of defence is every action upon it. Each instance of forced action must answer a question it is not accustomed to being asked: am I repairing a system genuinely deformed into negative curvature, or am I picking at a wound that would have healed? The modern conviction that a relation ``must be worked at,'' that maintenance is the default and letting-go the deviation, has the burden exactly backwards.
\end{claim}

This is, in fact, the deepest charge the paper lays against the logic of capital and of instrumental reason, and it is worth stating as such. That logic cannot believe that any system will come right of itself; it is constitutionally unable to grant the ground state its positive curvature. And so, by ceaseless management, optimization, and control, it bends every cycle that might have healed into one that now genuinely requires its perpetual management---manufacturing the very dependency it then presents itself as remedying. The relation that ``must be worked at'' becomes, under enough working-at, a relation that can no longer hold itself without the work; the managed system forgets how to float. The reversal of the burden of proof is thus not only an epistemic correction but a diagnosis of how the anxious, managing disposition produces the helplessness it claims to address---the same structure, at the scale of a whole rationality, as the anxious hand that rots the rose's root.

\subsection{Why there is no metric: the honest limit}
\label{sec:no-metric}

But the reversal must not become a licence, and here the section turns sharply upon itself, to forestall an over-strong claim that an earlier formulation of this very argument committed and that must now be retracted. It is tempting to convert the foregoing into a diagnostic rule: that the \emph{amount of forced action a cycle requires to sustain itself} is a reliable marker of its curvature---the good cycle (incentive-compatible, positive-curvature) tending toward self-maintenance with little forcing, the vicious cycle (incentive-incompatible, negative-curvature) demanding ceaseless forcing to prop up its counterfeit closure. There is something to this, and the commercial instance of \xref{sec:praxis} will draw on it; but as a general rule it overreaches, and the overreach must be retracted before it does damage.

The trouble is that one and the same external phenomenon---a high cost of maintenance---answers to at least three quite different internal states, and the external observation cannot, of itself, distinguish them.

\begin{claim}[The maintenance cost is a fallible marker, never a criterion]
The forced action a cycle requires to sustain itself is, at most, an external \emph{marker} of the sign of its curvature, and never a \emph{criterion} of it. One and the same high maintenance cost may answer to three quite different internal states. It may be the building of positive curvature in a founding period---a newly planted rose does need tending, and a young bond does need work, without this impugning its sustainability. It may be the repair of a positive-curvature system under external shock---a sound bond struck by illness or crisis needs upholding, and the upholding is no mark of viciousness. Or it may be the masking of an endogenous negative curvature---the vicious cycle forcing ceaselessly to hide its breach. The external phenomenon \emph{underdetermines} the internal property; the causal inference from marker to property is not deductively valid, and is, in the human case, scarcely susceptible of proof. The maintenance cost is therefore an invitation to diagnose, never the diagnosis itself.
\end{claim}

This limit, properly understood, is not a flaw in the framework but a corollary of its deepest commitment---and the section now makes the corollary thematic, for it governs the whole of the practical doctrine to come. The geometric framework offers a structural \emph{characterization} of what the good circle \emph{is} (a positive holonomy, an endogenous spiral); it does not, and cannot, offer an operational \emph{protocol} for measuring that property from outside.

\begin{claim}[The unobservability of curvature is a demand of the praxis, not a defect of the theory]
If the good and the vicious circle could be told apart by a single external indicator, then ``diagnostic attention''---the sustained, fallible, first-personal labour of reading the curvature---would be superfluous; one could simply consult the metric. But it is precisely the irreplaceability of this attention that will constitute the core of the practice of sustainability. A criterion outsourceable to ``check the maintenance cost'' would itself violate the paper's central claim, for it would settle, by an external sign, what can only be read along a lived path---it would be, at the level of method, a settling failure. The unobservability of the curvature is therefore not an embarrassment the theory must apologize for but a consequence the theory requires: there is no shortcut of judgement, and the sign of the curvature can be read only in a sustained, fallible, first-personal attention, exercised again and again. To mistake the structural ``what it is'' for an operational ``how to measure'' is the positivist's temptation, and the framework's refusal of a metric is its refusal of that temptation.
\end{claim}

So the epistemology arrives, by its own route, at a result that binds the practice to follow, and it is worth drawing the binding explicitly, for it dissolves the usual order of theory and practice. Knowing the curvature is not a measurement but an attention; and because it is an attention---sustained, fallible, first-personal, unoutsourceable---it is not a preliminary to the practice but already the practice itself. \emb{The epistemology is not the prelude to the practice; it is the content of the practice.} To sustain a relation is, in the end, to keep reading its curvature---which is to say, to attend; and the later sections on perception (\xref{sec:perception}) and on the disposition of \emph{wu wei} (\xref{sec:praxis}) will be, in large part, the unfolding of what this attending concretely involves. The attention that the epistemology here identifies as the only access to the curvature is the same attention that the semiotics will recognize as the witness the poem demands, and that the series' earlier papers described as joint attention. There is, in this convergence, a first intimation of the paper's recurring discovery: that to know a relation rightly and to sustain it rightly are not two acts but one, and that both are forms of an attention that declines to possess what it attends to.
